\rok{Јануар 2024.}{G}

Други практични тест из Алгоритама и структура података

Други практични тест одржан 12.01.2024.

\zadatak{1}{Конверзија листе у матрицу суседства}
Написати алгоритам који ће вршити директну конверзију листе суседства у матрицу суседства. Алгоритам као улазни параметар треба да прихвати динамички низ чворова, а треба да резултује креираном матрицом.

\textbf{Додатне напомене за решавање задатка:}
\begin{itemize}
	\item Написати алгоритам.
	\item У Main методи креирати произвољан граф. Обезбедити начин за тестирање и проверу нове функционалности, односно позив \texttt{convertListToMatrix} методе.
	\item У template-у је закоментарисана метода \texttt{convertListToMatrix} која треба да садржи имплементацију конверзије листе у матрицу суседства.
\end{itemize}



\zadatak{2}{Исписивање свих парних вредности чворова}
У оквиру стандардне структуре класе \texttt{AVL} написати имплементацију методе која ће омогућити проналажење и исписивање чворова са парном вредношћу.

\textbf{Додатне напомене за решавање задатка:}
\begin{itemize}
	\item У оквиру класе \texttt{AVL} креирати јавну методу са потписом:
	\begin{Verbatim}[fontsize=\footnotesize, numbers=left, xleftmargin=10mm]
public void printEvenNodes()
	\end{Verbatim}
	\item Из претходно креиране јавне методе треба да се позива приватна метода са потписом:
	\begin{Verbatim}[fontsize=\footnotesize, numbers=left, xleftmargin=10mm]
private Node printEvenNodesAVL(Node curr)
	\end{Verbatim}
	\item Приватна метода треба да садржи имплементацију алгоритма којим ће се омогућити пролазак кроз стабло од корена ка листовима и њихово исписивање уколико им је парна вредност.
	\item Листом се класификује сваки чвор који нема ниједно дете.
	\item Може се креирати помоћна променљива на нивоу класе \texttt{AVL} која ће помоћи у пребројавању.
	\item Алгоритам \textbf{ОБАВЕЗНО} писати у оквиру класе \texttt{AVL}.
\end{itemize}

\textbf{Примери:}

\begin{center}
\begin{minipage}{\textwidth}
\centering
\begin{forest}
for tree={
	circle,
	draw,
	thick,
	fill=gray!20,
	minimum size=8mm,
	inner sep=1pt,
	s sep=8mm,
	l sep=12mm,
	edge={-, thick},
	font=\small
}
[8
	[4
		[2
			[, phantom]
			[3]
		]
		[5]
	]
	[9
		[, phantom]
		[11]
	]
]
\end{forest}

\vspace{0.3cm}
\textbf{Слика 1.} Пример AVL-а.
\end{minipage}
\end{center}

\begin{itemize}
	\item \textbf{Улаз 1:} AVL са слике 1
	\item \textbf{Излаз 1:} 8, 4, 2
\end{itemize}

\begin{center}
\begin{minipage}{\textwidth}
\centering
\begin{forest}
for tree={
	circle,
	draw,
	thick,
	fill=gray!20,
	minimum size=8mm,
	inner sep=1pt,
	s sep=8mm,
	l sep=12mm,
	edge={-, thick},
	font=\small
}
[10
	[4
		[2
			[3]
			[, phantom]
		]
		[6
			[5]
			[9]
		]
	]
	[20
		[12
			[, phantom]
			[14]
		]
		[29
			[27]
			[, phantom]
		]
	]
]
\end{forest}

\vspace{0.3cm}
\textbf{Слика 2.} Пример AVL-а.
\end{minipage}
\end{center}

\begin{itemize}
	\item \textbf{Улаз 2:} AVL са слике 2
	\item \textbf{Излаз 2:} 10, 4, 2, 6, 20, 12, 14
\end{itemize}



\zadatak{3}{Укрштање секвенци са очувањем подсеквенце}
Написати алгоритам који укршта две секвенце, уколико је то могуће. Приликом укрштања неопходно је у обе секвенце очувати одређену подсеквенцу, па је њеном позицијом одређена позиција укрштања. Услов за укрштање секвенци је да је естимирана сличност секвенци које се укрштају мања од 0,9. У конзоли исписати укрштене секвенце или поруку да укрштање није могуће извести уколико је естимирана сличност секвенци које се укрштају већа од 0,9. Алгоритам имплементирати уз претпоставку да дата subsekventa сигурно постоји у секвенцама које су дате и појављује се само једном у њима.

\textbf{Додатне напомене за решавање задатка:}
\begin{itemize}
	\item Написати алгоритам.
	\item У Main методи обезбедити тестирање, позивом методе са параметрима који су дати у примерима.
	\item Користити \texttt{search} методу класе \texttt{KarpRabinAlgorithm} и класе \texttt{MinHash} и \texttt{Jaccard}, које су дате у шаблону.
	\item Имплементацију писати у оквиру закоментарисане методе:
	\begin{Verbatim}[fontsize=\footnotesize, numbers=left, xleftmargin=10mm]
public static void intersectSequencesAndKeepSubsequence(string str1, string str2, string substr)
	\end{Verbatim}
	\item Приликом креирања \texttt{MinHash} објекта, за вредности параметара \texttt{universeSize} и \texttt{numHashFunctions} прослеђивати вредности 12 и 4.
	\item Параметар \texttt{q} функције \texttt{search} поставити на \texttt{substr.Length+1}.
	\item Дозвољено је изменити повратни тип методе \texttt{search} (нпр. уместо \texttt{void} имати \texttt{List<int>}).
\end{itemize}

\textbf{Напомене које могу бити од помоћи приликом решавања задатка:}
\begin{itemize}
	\item За конверзију string-а у низ карактера користити \texttt{ToCharArray()} по принципу: \texttt{str.ToCharArray()}.
	\item За конверзију низа карактера у string користити \texttt{new string(char[] arrayOfChars)} по принципу: \texttt{new string(charArray)}.
	\item Карактере није потребно експлицитно конвертовати у целобројни тип приликом убацивања у \texttt{List<uint>} структуру.
\end{itemize}

\textbf{Пример:}
\begin{itemize}
	\item \textbf{Улаз 1:} \texttt{"ABHCGTMW", "DHMWCGTR", "CGT"}
	\item \textbf{Излаз 1:} \texttt{Ukrštene sekvence su: ABHCGTR i DHMWCGTMW}
	\item \textbf{Образложење:} Секвенце \texttt{ABHCGTMW} и \texttt{DHMWCGTR} имају естимирану сличност мању од 0,9. Место укрштања одређено је крајњом позицијом подсеквенце \texttt{CGT}.
	\item \textbf{Улаз 2:} \texttt{"ABHCGTMW", "ABHWCGTM", "CGT"}
	\item \textbf{Излаз 2:} \texttt{Nije moguće izvršiti ukrštanje.}
	\item \textbf{Образложење:} Сличност секвенци \texttt{ABHCGTMW} и \texttt{ABHWCGTM} је преко 0,9.
\end{itemize}



\sablon{Шаблон другог теста јануар 2024. група G}{2023/Januar/GrupaG/AISP2023_Test2_GrupaG.linq}
